\section{数据与表示里被悄悄用掉的假设}
\label{sec:18-Synthesis/01-System-Lifecycle/02}

信贷风控团队拿到了三年的历史放款数据，三十万条，每条都带着真实的还款结果。新模型在留出集上把坏账识别的 AUC 做到了 \(0.83\)，比线上那套用了五年的规则引擎高出一大截。上线半年，实际坏账率几乎没动。

问题不在模型。问题在于这三十万条记录里，每一条都是被旧规则批准过的申请。被旧规则拒掉的那些人从来没有产生过还款结果，因此也从来没有进入训练集。模型学到的是 \(P(\text{违约}\given \text{特征},\ \text{旧规则通过})\)，而上线之后要它回答的是 \(P(\text{违约}\given \text{特征})\)。这两个条件分布可以差得很远，尤其是在旧规则拒得最狠的那些区域——恰恰是新模型最需要有判断力的区域。

数据看上去是这条生命周期里最客观的一环：它是记录下来的事实，不像目标函数那样明显带着人的选择。这个印象是错的。从真实世界到模型输入之间隔着四层，每一层都用掉了一个没有被写进文档的假设。

\begin{center}
\begin{tikzpicture}[x=1cm,y=1cm,font=\small,>=stealth]
  \draw[black!45,fill=blue!8] (0,3.3) rectangle (4.2,4.3);
  \node[align=center] at (2.1,3.8) {真实世界中的总体};
  \draw[black!45,fill=blue!8] (0,1.95) rectangle (4.2,2.95);
  \node[align=center] at (2.1,2.45) {被记录下来的那部分};
  \draw[black!45,fill=blue!8] (0,0.6) rectangle (4.2,1.6);
  \node[align=center] at (2.1,1.1) {数据表中的行与列};
  \draw[black!45,fill=orange!10] (0,-0.75) rectangle (4.2,0.25);
  \node[align=center] at (2.1,-0.25) {模型看到的向量空间};
  \foreach \y in {3.3,1.95,0.6} \draw[->,black!55] (2.1,\y) -- (2.1,\y-0.35);
  \node[align=left,anchor=west,font=\footnotesize,black!70] at (4.7,3.12)
    {采样机制：谁进入了数据，谁永远不会};
  \node[align=left,anchor=west,font=\footnotesize,black!70] at (4.7,1.77)
    {编码与填充：假设缺失是随机的、类别可离散};
  \node[align=left,anchor=west,font=\footnotesize,black!70] at (4.7,0.42)
    {标准化与选择：这一步用了哪些行的信息};
  \node[align=left,anchor=west,font=\footnotesize,black!70] at (4.7,-0.25)
    {基与度量：``相近''在这里到底是什么意思};
\end{tikzpicture}

\emph{每一支向下的箭都用掉了一个没被写下来的假设。数据的客观性只到最上面那一层为止，往下每一步都是选择。}
\end{center}

\subsection{采样机制决定了你手上的分布是什么}

上图第一支箭是最致命的，因为它造成的偏差在数据内部完全看不出来。三十万条记录彼此一致、没有异常值、字段完整——被拒掉的那二十万人不会在任何地方留下痕迹，甚至不会体现为缺失值。

\hyperref[sec:05-Probability/01-Probability-Spaces/03]{``从现实问题到概率模型''}给出的第一个动作是把``什么可能发生''说清楚，也就是先确定样本空间再谈概率。这个动作在工程里常被跳过，因为大家默认``数据表就是样本空间''。而数据表其实是样本空间与一套记录机制的复合物。\hyperref[sec:11-Mathematical-Statistics/01-Statistical-Models/01]{``数据来自怎样的生成机制''}把这件事写成统计模型的第一句话：任何推断都是相对于一个假定的生成机制而言的，换了机制，同一份数据支持的结论就变了。

选择性记录有很多张面孔，但机制是同一个。二战期间统计学家看返航飞机的弹孔分布，最该加装甲的地方恰恰是弹孔最少的部位——因为打在那里的飞机没有回来。产品团队分析``活跃用户为什么留存高''，得到的结论对已经流失的人没有任何解释力。医院数据里两种互不相关的疾病显出负相关，仅仅因为只有病得足够重的人才会住院。这些例子里没有一个是数据质量问题，它们都是采样机制在起作用，而修复的唯一办法是把机制本身建模进来，或者去拿到被排除掉的那部分数据（风控里的做法是拿出一小部分额度做随机放款，用可控的成本买回缺失的那块分布）。

\subsection{预处理不是中性操作，它携带信息}

第三支箭上的那句话——``这一步用了哪些行的信息''——是离线与线上系统性脱节的最常见原因。

一个团队做完模型，离线 AUC \(0.88\)，上线 \(0.74\)。排查一圈没有分布漂移，没有代码缺陷。最后定位到标准化：均值与方差是在全量数据上算的，然后才做训练与验证的划分。验证折里每一个样本都对那两个统计量有过贡献，于是模型在验证时看到的输入已经含有验证集自身的信息。把标准化搬进每一折内部重做，离线指标掉到 \(0.76\)——和线上对上了。掉下去的那 \(0.12\) 从来就不属于模型。

标准化只是最容易发现的一种。缺失值用全量中位数填充、特征按与全量标签的相关性筛选、超参在包含验证折的数据上搜索、目标编码用了全量的类别均值、时间序列按随机而非按时间划分，全都是同一件事的不同写法：\textbf{任何在划分之前动过全量数据的步骤，都把``未见''这个前提破坏掉了}。\hyperref[sec:11-Mathematical-Statistics/06-Resampling-and-Model-Selection/03]{``交叉验证评估的是完整训练流程''}把这条讲成一句判断——被评估的对象是那条从原始数据到预测的完整流程，不是流程末端的那个模型对象。只要有任何一个环节站在折外，交叉验证给出的就是一个系统性乐观的数。

有一个便宜的自检可以抓住大部分泄漏：把标签随机打乱，整条流程原样重跑一遍，指标应当回到随机水平。回不去就说明流程里有通道让标签信息绕过了模型。这个检查花不了多少算力，却能挡掉这一步最贵的几种错误。

\subsection{表示的选择就是度量的选择}

数据表变成向量之后，事情还没完。同一份数据用不同的基、不同的量纲，``两个样本相近''这句话的含义完全不同——而下游几乎所有算法（最近邻、聚类、核方法、正则化）都建立在这句话上。

\hyperref[sec:03-Linear-Algebra/02-Vector-Spaces/04]{``线性变换与换基：对象不变，坐标可以改变''}讲的是同一个对象在不同坐标下写出来不一样，但对象本身没变。这条在数据里要反过来用：坐标变了，对象确实没变，可\emph{距离变了}。把``年收入''按元和按万元记录，欧氏距离里它的权重差四个数量级，于是标准化不是清洗步骤而是一次度量选择——它默认每个特征同等重要。同理，\hyperref[sec:03-Linear-Algebra/03-Orthogonality/02]{``投影与最近点：从一条直线走向一般子空间''}给出的判断是最近点由内积决定；换一个内积，``最近''就指向另一个点。特征工程做的事本质上就是选内积，只是很少有人这样描述它。

独立性同样不是数据的固有属性。\hyperref[sec:05-Probability/05-Joint-Distributions/02]{``随机变量的独立与依赖''}里那条容易被忽略的判断是：独立性是相对于一组变量和一个分布而言的，加一个条件变量、换一个子群体，原本独立的两个量可能立刻相关起来。所以``这两个特征不相关，可以分开处理''这种说法，只在当前这份数据的当前切法下成立。

维度一高，直觉开始整体失效。在高维各向同性的数据里，任意两点之间的距离会向同一个值集中：最近邻与最远邻的距离之比趋于 \(1\)，于是``最近的那个''不再携带多少信息。这不是算法实现的问题，是几何本身的性质。真实数据之所以还能用最近邻，靠的是它并不真的填满高维空间——\hyperref[sec:12-Real-Analysis-Support/10-Essential-Structures/02]{``低维结构为什么普遍存在''}给出的理由是：数据由少数几个自由度生成，因此集中在一个远低于名义维数的子集附近。\hyperref[sec:13-Machine-Learning/04-Dimensionality-Reduction/04]{``流形假设与非线性降维''}把这句话写成可用的形式，同时诚实地指出它是一个假设而不是定理，在某些数据上就是不成立。

有一件事必须与低维结构分开记：即使数据真的低维，你从样本里\emph{看到}的谱也未必是真实的谱。\hyperref[sec:11-Mathematical-Statistics/09-Random-Matrix-Theory/01]{``样本协方差在高维为何会失真''}给出的结论是，当特征数与样本量可比时，纯噪声数据的样本协方差矩阵的特征值也会散布在一个很宽的区间上，最大的那几个显著大于 \(1\)。这意味着碎石图上那个漂亮的肘部可能完全是噪声造出来的，跟着它选主成分个数会一次次选出并不存在的结构。

\subsection{这一步做错时，故障长什么样}

离线与线上系统性差一截，而且是稳定地差同一个幅度——这几乎总是泄漏，不是漂移。漂移的特征是差距随时间变化，泄漏的特征是差距从第一天起就在那里。开头那个 \(0.88\) 与 \(0.74\) 是教科书式的形态。

换一批数据重训之后表现大幅波动，说明模型依赖的信号里有相当一部分是这批数据特有的。可能是采样机制在两批数据之间变了，也可能是特征选择每次都在噪声里挑出不同的一组特征。后者在特征数远大于样本数时几乎必然发生。

降维之后的可视化图上出现清晰的簇，但业务上找不到对应的含义；或者主成分个数每换一批数据就变一次。这是谱失真的典型症状，值得先做一次零假设检查：把每列特征各自随机打乱以破坏列间相关性，重新算一遍谱，看看那个``肘部''是不是照样出现。

最后一种最难查：模型在整体上表现良好，在某个子群体上明显更差。这通常意味着那个子群体在训练数据里的比例与真实比例不同，或者它的标签是通过另一条路径产生的。整体指标对这种问题完全不敏感，只有分层看才看得见。

数据与表示落定之后，下一个决定是选哪一族模型。这个决定听上去是关于算力和精度的，实际上是关于世界形状的：每一个模型族都在替你断言数据应当长成什么样。
